%% kapa: exam=1 round=26 year=2015 period=1 form=A subject=economics subject_ko=경제원론 seq=38 number=78
%% source: Q-Net 감정평가사 기출문제 (cst00309) artlSeq=5138001 file=extracted/1008/5138001/제26회_감정평가사 1차_기출/제26회_감정평가사 1차_기출/1교시_감평_A형.pdf
%% conversion: kordoc 4.24.51 -> tools/gen1.py (auto); review=proofread-sonnet; flags=none
\begin{question}{78}
甲국과 乙국의 무역 개시 이전의 $X$재와 $Y$재에 대한 단위당 생산비가 다음과 같다. 무역을 개시하여 두 나라 모두 이익을 얻을 수 있는 교역조건 $(P_{X}/P_{Y})$에 해당하는 것은? (단, $P_{X}$는 $X$재의 가격이고, $P_{Y}$는 $Y$재의 가격이다.)
\par{\small\centering\begin{tabular}{|c|c|c|}\hline
 & $X$재 & $Y$재 \\ \hline
甲국 & 5 & 10 \\ \hline
乙국 & 8 & 13 \\ \hline
\end{tabular}\par}
\begin{choices}
  \item 0.45
  \item 0.55
  \item 0.65
  \item 0.75
  \item 0.85
\end{choices}
\end{question}
\answer{2}
